\documentclass[11pt]{article}
\usepackage[a4paper,margin=24mm]{geometry}
\usepackage{amsmath,amssymb,amsthm,mathtools,bm}
\usepackage{booktabs,longtable,array}
\usepackage{enumitem}
\usepackage{xcolor}
\usepackage{microtype}
\usepackage[hidelinks]{hyperref}

\setlength{\emergencystretch}{3em}
\newtheorem{proposition}{Proposition}[section]
\newtheorem{lemma}[proposition]{Lemma}
\newtheorem{corollary}[proposition]{Corollary}
\theoremstyle{definition}
\newtheorem{definition}[proposition]{Definition}
\newtheorem{gate}[proposition]{Fail-closed gate}
\theoremstyle{remark}
\newtheorem{remark}[proposition]{Remark}

\newcommand{\dd}{\mathrm d}
\newcommand{\ii}{\mathrm i}
\newcommand{\Tr}{\operatorname{Tr}}
\newcommand{\cE}{\mathcal E}
\newcommand{\cL}{\mathcal L}
\newcommand{\cH}{\mathcal H}
\newcommand{\cU}{\mathcal U}
\newcommand{\RR}{\mathbb R}
\newcommand{\ZZ}{\mathbb Z}
\newcolumntype{L}[1]{>{\raggedright\arraybackslash}p{#1}}

\title{AP-E3 Nonlinear Follow-up:\\
A Charged Two-Colour EFT Proxy, Diquark Gap,\\
and the Necessary \(B=1\) Soliton-Hessian Gate}
\author{Route F research note}
\date{16 July 2026}

\begin{document}
\maketitle

\begin{abstract}
We construct and solve a deliberately limited low-energy proxy for the
unsettled strong-phase and soliton-stability gates of the charged two-colour
AP-E3 candidate.  A real \(O(4)\) meson multiplet is supplemented by a
complex \(U(1)_g\)-charge-two diquark.  Its homogeneous potential yields an
analytic phase boundary, vacuum condensates, and the exact tree-level gaps
\[
 m_\pi^2=\frac{h}{f},\qquad
 m_\sigma^2=m_\pi^2+2\lambda f^2,\qquad
 m_\Delta^2=m_\pi^2+\mu_g^2.
\]
In the mesonic phase we take the nonlinear \(SU(2)\) limit, add the Skyrme
term, solve the full \(B=1\) hedgehog boundary-value problem, verify baryon
number and Derrick's virial identity, and diagonalise two quadratic blocks:
the generalized spherically symmetric meson Hessian and the most dangerous
charged-diquark channel.  Grid and volume sequences converge; an intentionally
over-attractive diquark coupling produces a negative eigenvalue and therefore
tests that the audit can fail.

The term nonperturbative is used only in the narrow classical sense that no
small-amplitude or small-profile expansion is used in solving the nonlinear
boundary-value problem.  This is not lattice two-colour QCD, does not
integrate the microscopic gauge theory, and does not close the full
three-dimensional coupled Hessian.  Moreover, the product target used here
excludes the global \(S^5\) escape direction by assumption.  The result is
consequently a useful necessary gate and a reproducible falsifier, not a proof
of the microscopic charged phase or of a stable quantum baryon.
\end{abstract}

\tableofcontents

\section{Epistemic boundary and questions answered}

The anomaly-consistent charged two-colour construction fixes ultraviolet
quantum numbers but not the infrared vacuum.  Gauge-anomaly cancellation
does not prove a mesonic condensate, a positive diquark gap, or a stable
\(B=1\) saddle.  Effective descriptions of two-colour dynamics explicitly
retain meson and diquark directions
\cite{SuenagaEtAl2023Proxy,SuenagaEtAl2024Proxy}; the pseudoreal theory is
therefore not equivalent to ordinary three-colour QCD.  This note asks three
narrower questions.
\begin{enumerate}[label=(Q\arabic*)]
\item Can one exhibit a charged effective potential with a transparent
  meson--diquark phase boundary and a positive benchmark Hessian?
\item Does its declared mesonic nonlinear subsector possess a numerically
  converged \(B=1\) hedgehog satisfying the exact topology and scale tests?
\item Are its radial meson and charged-diquark quadratic blocks locally
  nonnegative, and can the same implementation detect an unstable control?
\end{enumerate}

The answers are yes inside the proxy.  The following statements remain
false: a lattice or continuum nonperturbative solution of the microscopic
\(SU(2)_c\times U(1)_g\) gauge theory; a UV derivation of every effective
coefficient; a full three-dimensional Hessian; a fermion-determinant audit;
and exclusion of finite-amplitude unwinding paths.

\begin{gate}[Promotion boundary]
A positive result below sets only the nonlinear-EFT and radial-Hessian
necessary gates to true.  It leaves the full three-dimensional Hessian,
lattice QC2D, global unwinding, and physics-promotion gates false.
\end{gate}

\section{Charged homogeneous EFT proxy}

\subsection{Fields, charge, and potential}

Let
\[
 \varphi=(\sigma,\pi_1,\pi_2,\pi_3,d_1,d_2),\qquad
 \Delta=\frac{d_1+\ii d_2}{\sqrt2}.
\]
The mesons are \(U(1)_g\)-neutral and \(\Delta\) has charge
\(q_\Delta=2\), as expected for a two-quark operator.  Thus
\[
 D_\mu\Delta=(\partial_\mu-2\ii e_g b_\mu)\Delta.
\]
With canonical real-field normalization, the proxy Lagrangian is
\begin{equation}
 \cL_{\rm hom}
 =\frac12(\partial_\mu\sigma)^2+\frac12(\partial_\mu\bm\pi)^2
 +|D_\mu\Delta|^2-\frac14F_{\mu\nu}F^{\mu\nu}-V,
 \label{eq:hom-lagrangian}
\end{equation}
\begin{equation}
 V=\frac{\lambda}{4}
 \left(\sigma^2+\bm\pi^2+2|\Delta|^2-v^2\right)^2
 -h\sigma+\mu_g^2|\Delta|^2.
 \label{eq:potential}
\end{equation}
Here \(\mu_g^2\) is a matched Pauli--G\"ursey-breaking coefficient.  Its
sign and magnitude are not inferred merely from the gauge charge.
Equation~\eqref{eq:potential} is a diagnostic model, not a claimed exact
chiral Lagrangian of the UV candidate.

\subsection{Mesonic stationary point and exact vacuum Hessian}

At \(\bm\pi=\Delta=0\), choose the positive root \(f\) of
\begin{equation}
 \lambda(f^2-v^2)f=h.
 \label{eq:f-equation}
\end{equation}
Define \(m_\pi^2=h/f=\lambda(f^2-v^2)\).  Direct differentiation of
Eq.~\eqref{eq:potential} gives
\begin{equation}
 \left.\nabla^2V\right|_{(f,\bm0,0)}
 =\operatorname{diag}\left(
 m_\sigma^2,
 m_\pi^2,m_\pi^2,m_\pi^2,
 m_\Delta^2,m_\Delta^2\right),
 \label{eq:vacuum-hessian}
\end{equation}
where
\begin{equation}
 m_\sigma^2=m_\pi^2+2\lambda f^2,\qquad
 m_\Delta^2=m_\pi^2+\mu_g^2.
 \label{eq:gaps}
\end{equation}
Therefore the charged direction is locally stable precisely when
\begin{equation}
 \mu_g^2>-m_\pi^2.
 \label{eq:local-phase-boundary}
\end{equation}
The two identical diquark eigenvalues follow from unbroken \(U(1)_g\), and
linear meson--diquark mixing is forbidden by charge.

\subsection{Condensed branch and continuous onset}

When \(\Delta\ne0\), the stationary equations read
\begin{equation}
 \lambda(R^2-v^2)\sigma=h,\qquad
 [\lambda(R^2-v^2)+\mu_g^2]d_a=0,
 \quad R^2=\sigma^2+d_1^2+d_2^2.
\end{equation}
For \(\mu_g^2<0\), the charged branch is consequently
\begin{equation}
 \sigma=-\frac{h}{\mu_g^2},\qquad
 d_1^2+d_2^2=v^2-\frac{\mu_g^2}{\lambda}
 -\frac{h^2}{(\mu_g^2)^2}.
 \label{eq:condensed-branch}
\end{equation}
The right-hand side vanishes continuously at
\(\mu_g^2=-h/f=-m_\pi^2\).  Below that point \(U(1)_g\) is Higgsed and the
homogeneous Hessian contains the corresponding angular zero mode before
gauge fixing.  This analytic onset is used as a positive and negative
classification test in the scan.

\subsection{Explicit
\texorpdfstring{\((e_g,v_g/\Lambda_c)\)}{(eg, vg/Lambda-c)}
sensitivity chart}

The EFT coefficient \(\mu_g^2\) is not yet known from the charged microscopic
theory.  To avoid hiding that ignorance while still scanning the requested
coordinates, we declare the matching ansatz
\begin{equation}
 \frac{\mu_g^2}{\Lambda_c^2}
 =\frac{\mu_{\rm PG,0}^2}{\Lambda_c^2}
 +c_g\left(e_g\frac{v_g}{\Lambda_c}\right)^2.
 \label{eq:matching-ansatz}
\end{equation}
The numerical chart uses
\(\mu_{\rm PG,0}^2/\Lambda_c^2=-0.40\), \(c_g=0.24\), and
\(\Lambda_c=f=1\).  This makes both sides of the analytic phase boundary
visible.  It is a sensitivity coordinate chart, not a loop matching result.
Every cell records \(e_g\), \(v_g/\Lambda_c\), \(\langle\sigma\rangle\),
\(|\langle d\rangle|\), and the smallest homogeneous Hessian eigenvalue.
Table~\ref{tab:charged-chart} shows the complete deterministic chart.
Here D denotes the diquark-condensed branch and M the mesonic branch.
The zero on D is the ungauge-fixed \(U(1)_g\) angular mode, not a negative
curvature.

\begin{longtable}{@{}rrrrrr@{}}
\caption{Declared charged-coordinate sensitivity chart.}
\label{tab:charged-chart}\\
\toprule
\(e_g\) & \(v_g/\Lambda_c\) & \(\mu_g^2/\Lambda_c^2\) &
phase & \(\sigma/\Lambda_c\) & \((|d|,\lambda_{\min})\)\\
\midrule
\endfirsthead
\toprule
\(e_g\) & \(v_g/\Lambda_c\) & \(\mu_g^2/\Lambda_c^2\) &
phase & \(\sigma/\Lambda_c\) & \((|d|,\lambda_{\min})\)\\
\midrule
\endhead
0.00 & 0.50 & -0.4000 & D & 0.625000 & (0.796477,0)\\
0.00 & 1.00 & -0.4000 & D & 0.625000 & (0.796477,0)\\
0.00 & 2.00 & -0.4000 & D & 0.625000 & (0.796477,0)\\
0.00 & 3.00 & -0.4000 & D & 0.625000 & (0.796477,0)\\
0.25 & 0.50 & -0.3963 & D & 0.630915 & (0.791405,0)\\
0.25 & 1.00 & -0.3850 & D & 0.649351 & (0.775141,0)\\
0.25 & 2.00 & -0.3400 & D & 0.735294 & (0.688725,0)\\
0.25 & 3.00 & -0.2650 & D & 0.943396 & (0.335416,0)\\
0.50 & 0.50 & -0.3850 & D & 0.649351 & (0.775141,0)\\
0.50 & 1.00 & -0.3400 & D & 0.735294 & (0.688725,0)\\
0.50 & 2.00 & -0.1600 & M & 1.000000 & (0,0.090000)\\
0.50 & 3.00 &  0.1400 & M & 1.000000 & (0,0.250000)\\
0.75 & 0.50 & -0.3663 & D & 0.682594 & (0.743936,0)\\
0.75 & 1.00 & -0.2650 & D & 0.943396 & (0.335416,0)\\
0.75 & 2.00 &  0.1400 & M & 1.000000 & (0,0.250000)\\
0.75 & 3.00 &  0.8150 & M & 1.000000 & (0,0.250000)\\
1.00 & 0.50 & -0.3400 & D & 0.735294 & (0.688725,0)\\
1.00 & 1.00 & -0.1600 & M & 1.000000 & (0,0.090000)\\
1.00 & 2.00 &  0.5600 & M & 1.000000 & (0,0.250000)\\
1.00 & 3.00 &  1.7600 & M & 1.000000 & (0,0.250000)\\
\bottomrule
\end{longtable}

\section{Mesonic nonlinear limit and the
\texorpdfstring{\(B=1\)}{B=1} saddle}

\subsection{Why a product-target proxy is used}

In the positive-gap phase, freeze the heavy radial mode and set
\(\Delta=0\).  Write
\begin{equation}
 U=\frac{\sigma+\ii\bm\tau\cdot\bm\pi}{f}\in SU(2).
\end{equation}
We then use the standard sigma plus Skyrme energy
\cite{Skyrme1961,AdkinsNappiWitten1983}.  This is a product-target
construction: \(U\) remains an \(SU(2)\) coordinate when the spectator
\(\Delta\) fluctuates.  It therefore retains
\(\pi_3(SU(2))=\ZZ\).

That choice is useful for a local gate but logically costly.  In ungauged
two-colour QCD the Pauli--G\"ursey target for two flavours is
\(SU(4)/Sp(4)\simeq S^5\), for which \(\pi_3(S^5)=0\).  The present proxy
does not prove that the charged microscopic theory dynamically factorizes
its target, nor does it exclude a finite-amplitude path that restores an
\(S^5\)-like escape direction.  This is why local Hessian positivity is not
promoted to global soliton stability.

\subsection{Static energy and hedgehog reduction}

Take
\begin{equation}
 E_U=\int\dd^3x\left[
 \frac{f^2}{4}\Tr(\partial_iU\,\partial_iU^\dagger)
 +\frac{1}{32e^2}\Tr\!\left([L_i,L_j]^\dagger[L_i,L_j]\right)
 +\frac{f^2m_\pi^2}{2}\Tr(1-U)\right],
 \label{eq:skyrme-energy}
\end{equation}
where \(L_i=U^\dagger\partial_iU\).  For
\begin{equation}
 U(\bm x)=\cos F(r)+\ii\,\bm\tau\cdot\widehat{\bm x}\sin F(r),
 \qquad x=efr,\qquad \beta=\frac{m_\pi}{ef},
 \label{eq:hedgehog}
\end{equation}
the energy is
\begin{equation}
 E_U=\frac{4\pi f}{e}\int_0^\infty\dd x\,\cE_F,
 \label{eq:radial-energy}
\end{equation}
\begin{equation}
 \cE_F=
 \frac12x^2F'^2+\sin^2F(1+F'^2)
 +\frac{\sin^4F}{2x^2}
 +\beta^2x^2(1-\cos F).
 \label{eq:radial-density}
\end{equation}
The baryon number reduces to
\begin{equation}
 B=-\frac2\pi\int_0^\infty\sin^2F\,F'\,\dd x
 =\frac1\pi\left[F-\frac12\sin2F\right]_{x=\infty}^{x=0}.
 \label{eq:baryon-number}
\end{equation}
Thus \(F(0)=\pi\), \(F(\infty)=0\) imply \(B=1\).

\subsection{Euler--Lagrange equation and regular origin data}

Varying Eq.~\eqref{eq:radial-density} yields
\begin{equation}
 (x^2+2\sin^2F)F''
 +2xF'
 +2\sin F\cos F(F'^2-1)
 -\frac{2\sin^3F\cos F}{x^2}
 -\beta^2x^2\sin F=0.
 \label{eq:hedgehog-ode}
\end{equation}
A numerically stable origin condition cannot set \(F(\epsilon)=\pi\) while
allowing a finite nonzero slope.  Substitution of
\begin{equation}
 F(x)=\pi-a x+b x^3+O(x^5)
\end{equation}
into Eq.~\eqref{eq:hedgehog-ode} gives
\begin{equation}
 b=\frac{a(2a^4+4a^2+3\beta^2)}{30(1+2a^2)}.
 \label{eq:origin-series}
\end{equation}
The boundary-value solver treats \(a\) as an unknown parameter and imposes
both Eq.~\eqref{eq:origin-series} at \(x=\epsilon\) and \(F(L)=0\).

\subsection{Derrick identity and collapse negative control}

Split the dimensionless integral into \(E_2+E_4+E_0\), corresponding to
the two-derivative, four-derivative, and mass terms.  Under
\(F(x)\mapsto F(x/\lambda)\),
\begin{equation}
 E(\lambda)=\lambda E_2+\lambda^{-1}E_4+\lambda^3E_0.
\end{equation}
Any finite-size stationary solution must obey
\begin{equation}
 E_2-E_4+3E_0=0,\qquad
 E''(1)=2E_4+6E_0>0.
 \label{eq:derrick}
\end{equation}
Deleting the Skyrme contribution leaves
\(E'_{\rm no\;Skyrme}(1)=E_2+3E_0>0\), so shrinking lowers the energy.
The code checks both the virial identity and this deliberate collapse
control.

\section{Second variation and spectral problems}

\subsection{Spherically symmetric meson block}

Let \(F\mapsto F+\varepsilon\eta\) with
\(\eta(0)=\eta(L)=0\).  Write
\begin{equation}
 \cE_F=\frac12A(F,x)F'^2+\cU(F,x),\quad
 A=x^2+2\sin^2F,
\end{equation}
\begin{equation}
 \cU=\sin^2F+\frac{\sin^4F}{2x^2}
 +\beta^2x^2(1-\cos F).
\end{equation}
Twice differentiating the energy and integrating the mixed
\(\eta\eta'\) term by parts gives
\begin{equation}
 \delta^2E_F=\frac{4\pi f}{e}\int_0^L\dd x
 \left[A\eta'^2+V_F\eta^2\right],
 \label{eq:radial-second-variation}
\end{equation}
\begin{equation}
 V_F=\cU_{FF}+\frac12A_{FF}F'^2-\bigl(A_F F'\bigr)',
 \label{eq:VF-abstract}
\end{equation}
or, using the background equation,
\begin{align}
 V_F={}&2\cos2F
 +\frac{2}{x^2}\left(3\sin^2F\cos^2F-\sin^4F\right)
 +\beta^2x^2\cos F \notag\\
 &-2\cos2F\,F'^2-4\sin F\cos F\,F''.
 \label{eq:VF-explicit}
\end{align}
The same sigma and Skyrme terms give the radial kinetic norm
\(\int A\eta^2\dd x\).  Hence the displayed squared frequencies solve the
generalized Sturm--Liouville problem
\begin{equation}
 \cH_F\eta=\omega^2A\eta,\qquad
 \cH_F=-\frac{\dd}{\dd x}A\frac{\dd}{\dd x}+V_F.
 \label{eq:radial-generalized}
\end{equation}
At large \(x\), the continuum floor is \(\omega^2=\beta^2\).  A lowest
finite-box eigenvalue approaching \(\beta^2\) from above is therefore not a
localized breathing mode.

\subsection{Charged-diquark block}

The leading symmetry-allowed static spectator energy is chosen as
\begin{equation}
 E_\Delta^{(2)}
 =\frac{4\pi f}{e}\sum_{a=1}^2\frac12
 \int_0^\infty\dd x\,x^2
 \left[(d_a')^2+
 \{\alpha^2+\chi(1-\cos F)\}d_a^2\right],
 \label{eq:diquark-quadratic}
\end{equation}
where
\begin{equation}
 \alpha=\frac{m_\Delta}{ef}.
\end{equation}
The coefficient \(\chi\) measures the matched core interaction.  It is
independent input and is scanned with a tachyonic negative control.  Because
the background has \(\Delta=0\), the two real components are degenerate and
there is no quadratic meson--diquark mixing.

For angular momentum \(\ell\), set \(u_\ell=xd_\ell\).  The Hessian becomes
\begin{equation}
 \cH_{\Delta,\ell}
 =-\frac{\dd^2}{\dd x^2}+\alpha^2+\chi(1-\cos F)
 +\frac{\ell(\ell+1)}{x^2},
 \qquad u_\ell(0)=u_\ell(L)=0.
 \label{eq:diquark-operator}
\end{equation}
Thus \(\ell=0\) is the most dangerous scalar channel:
\(\cH_{\Delta,\ell}-\cH_{\Delta,0}\ge0\).  Positivity of its lowest
eigenvalue closes the charged-scalar block in this proxy.  It does
\emph{not} close nonradial meson, vector, gauge, radial-\(\sigma\), ghost, or
fermionic channels.  A positive bound mode satisfies
\begin{equation}
 0<\omega_{\Delta,0}^2<\alpha^2.
\end{equation}

\subsection{Gauge-field block at the declared saddle}

At \(\Delta=0\), the Maxwell field does not mix linearly with
\(d_1,d_2\).  With nonnegative Maxwell and optional Proca/Higgs quadratic
terms, its isolated tree-level block is nonnegative.  This observation does
not replace gauge fixing and ghost accounting in the full coupled
three-dimensional problem; those remain explicit blockers.

\section{Deterministic numerical experiment}

\subsection{Benchmark and algorithms}

The benchmark fixes
\begin{equation}
 f=e=1,\quad \lambda=6,\quad m_\pi=0.5,\quad
 \mu_g^2=0.56,\quad \chi=-2,
 \label{eq:benchmark}
\end{equation}
so that
\[
 v^2=\frac{23}{24},\quad h=\frac14,\quad
 (m_\pi,m_\sigma,m_\Delta)=(0.5,3.5,0.9),\quad
 (\beta,\alpha)=(0.5,0.9).
\]
The BVP is solved on \([\epsilon,L]\), \(\epsilon=10^{-4}\), using the
regular series and an adaptive collocation tolerance of \(3\times10^{-7}\).
The radial generalized problem uses linear finite elements with consistent
mass matrix.  The transformed diquark problem uses a centered
second-difference operator.  Sparse symmetric eigensolvers request the
algebraically lowest eigenvalues.

\subsection{Benchmark results}

The machine-readable output is authoritative for the final digits.  The
expected audit structure is
\begin{longtable}{@{}L{0.38\textwidth}L{0.25\textwidth}L{0.27\textwidth}@{}}
\toprule
Quantity & Result & Interpretation\\
\midrule
\endfirsthead
\toprule
Quantity & Result & Interpretation\\
\midrule
\endhead
Vacuum gaps & \((0.5,3.5,0.9)\) & positive meson, radial, diquark gaps\\
\(B\) and \(|B-1|\) &
\(0.999999999594,\ 4.06\times10^{-10}\) & topology/integration check\\
\(E_2-E_4+3E_0\) &
\(-1.55\times10^{-7}\) & relative residual \(2.52\times10^{-8}\)\\
lowest \(\omega_F^2\) & \(0.311400567\) & \(L=18,N=1200\) radial gate\\
lowest \(\omega_{\Delta,0}^2\) & \(0.696506101\) &
positive bound mode below \(\alpha^2=0.81\)\\
\(\chi=-5\) control & \(-1.430515810\) & charged tachyon detected\\
\bottomrule
\end{longtable}

Three grids with \(N=600,1200,2400\) at fixed \(L=18\) test discretization.
A second sequence with
\((L,N)=(14,700),(18,900),(22,1100)\) tests the box.  The soliton energy and
localized diquark eigenvalue must stabilize.  The first radial finite-box
state is instead expected to decrease toward the continuum floor
\(\beta^2=0.25\), remaining positive.

\begin{table}[htbp]
\centering
\begin{tabular}{@{}rrrr@{}}
\toprule
\(N\) & spacing & \(\omega_{F,\min}^2\) &
\(\omega_{\Delta,\min}^2\)\\
\midrule
600  & 0.02999983 & 0.311401592 & 0.696441268\\
1200 & 0.01499992 & 0.311400567 & 0.696506101\\
2400 & 0.00749996 & 0.311400310 & 0.696522301\\
\bottomrule
\end{tabular}
\caption{Fixed-box discretization sequence.  Both observed convergence orders
are approximately two; the machine-readable card records Richardson
extrapolates and error estimates.}
\label{tab:grid-convergence}
\end{table}

\begin{table}[htbp]
\centering
\begin{tabular}{@{}rrrrr@{}}
\toprule
\(L\) & \(N\) & \(\int\cE_F\dd x\) &
\(\omega_{F,\min}^2\) & \(\omega_{\Delta,\min}^2\)\\
\midrule
14 & 700  & 6.135914701 & 0.349739800 & 0.696547693\\
18 & 900  & 6.135914193 & 0.311400833 & 0.696489297\\
22 & 1100 & 6.135914184 & 0.291384645 & 0.696485381\\
\bottomrule
\end{tabular}
\caption{Fixed-spacing volume sequence.  The radial finite-box state moves
toward \(\beta^2=0.25\), whereas the localized charged mode stabilizes.}
\label{tab:volume-convergence}
\end{table}

\subsection{What the negative controls prove}

Three failure modes are built into the same calculation.
\begin{enumerate}
\item Setting \(\mu_g^2<-m_\pi^2\) makes the mesonic stationary point
  unstable and moves the homogeneous minimum to the charged branch.
\item Removing \(E_4\) violates the possibility of a finite-size Derrick
  stationary point and produces collapse.
\item Replacing the benchmark \(\chi=-2\) by \(\chi=-5\) makes the lowest
  charged-diquark eigenvalue negative.
\end{enumerate}
Consequently a green benchmark is not caused by an eigenvalue routine that
silently takes absolute values or by a phase classifier that always selects
the mesonic branch.

\section{Interpretation and remaining routes}

\subsection{What has been learned}

The proxy demonstrates a logically consistent local scenario:
\begin{equation}
\begin{gathered}
 m_\Delta^2>0,\quad \langle\Delta\rangle=0,\quad
 U(1)_g\ {\rm unbroken},\\
 B=1,\quad \delta^2E_F>0\ {\rm in\ the\ radial\ sample},\quad
 \delta^2E_\Delta>0\ {\rm for\ every\ scalar}\ \ell .
\end{gathered}
\end{equation}
It also quantifies how a sufficiently attractive soliton core can destroy
that scenario even while the homogeneous vacuum retains a positive diquark
gap.  Vacuum stability and soliton-core stability are therefore independent
gates.

\subsection{What has not been learned}

The following gaps cannot be repaired by increasing the radial grid.
\begin{enumerate}
\item The ansatz~\eqref{eq:matching-ansatz} and the coupling \(\chi\) need
  matching from the microscopic charged gauge theory.
\item The mesonic \(SU(2)\) product factor is assumed.  A faithful full target
  may allow global unwinding even when every local eigenvalue reported here
  is positive.
\item Nonradial meson deformations, the finite radial-\(\sigma\) field,
  vector mesons, gauge/ghost fields, and collective zero modes require a
  coupled three-dimensional Hessian.
\item Fermionic and gauge determinants can shift the classical spectrum.
\item A lattice calculation must measure the meson and diquark condensates,
  charged gap, baryon free energy, and finite-volume scaling.
\end{enumerate}

\subsection{Recommended next implementations}

There are three useful continuations, ordered by evidential strength.
\begin{enumerate}
\item \textbf{Lattice phase chart.}  Scan \((e_g,v_g/\Lambda_c)\) and measure
  \(\langle\bar qq\rangle\), a gauge-invariant dressed diquark correlator,
  the corresponding gap, and volume scaling.  This directly replaces the
  matching ansatz.
\item \textbf{Full EFT Hessian.}  Retain \(\sigma\), vector mesons and the
  massive \(U(1)_g\) field on a three-dimensional grid, gauge-fix, remove
  collective modes, and continue the lowest eigenpairs in \(\chi\).
\item \textbf{Global path search.}  Use a string or nudged-elastic-band method
  in the faithful field space between the \(B=1\) configuration and vacuum.
  This tests finite-amplitude escape that a Hessian cannot see.
\end{enumerate}

\section{Reproducibility ledger}

Run from the repository root:
\begin{verbatim}
python3 route_f/code/scan_ap_e3_charged_two_colour_proxy.py
\end{verbatim}
The program writes
\begin{verbatim}
route_f/output/ap_e3_charged_two_colour_proxy.json
route_f/output/ap_e3_charged_two_colour_proxy.md
\end{verbatim}
and records Python, NumPy, SciPy, source hashes, every mechanical assertion,
the phase chart, convergence sequences, Hessian spectra, negative controls,
and fail-closed booleans.

\bibliographystyle{unsrt}
\bibliography{route_f/tex/ap_e3_charged_two_colour_soliton_proxy}

\end{document}
